\section{Сценарии тестирования}

Для демонстрации работоспособности алгоритма рассмотрены сценарии, отражающие различные аспекты работы системы интернета вещей.

\subsection{Сценарий 1: Стабильная нагрузка (Baseline)}

\textbf{Описание}: В этом режиме интенсивность поступления задач постоянна ($\lambda = 2.5$), отказы узлов отсутствуют, а ресурсы узлов фиксированы. 

\textbf{Цель}: Оценить производительность алгоритма в оптимальных условиях.

\textbf{Ожидаемые результаты}:
\begin{itemize}
	\item Средняя задержка выполнения: $< 150$ мс
	\item Процент принятых задач: $> 90\%$
	\item Социальное благосостояние: растущий тренд
	\item Справедливость платежей: коэффициент Джини платежей $< 0.3$
\end{itemize}

\subsection{Сценарий 2: Динамическая нагрузка с пиками (Load Spikes)}

\textbf{Описание}: Нагрузка варьируется синусоидально и дополняется случайными всплесками спроса, моделирующие реальные события (срабатывание датчиков, запуск аналитики). 

\textbf{Цель}: Оценить адаптивность алгоритма при резком росте конкуренции за ресурсы.

\textbf{Ожидаемые результаты}:
\begin{itemize}
	\item Система должна адаптироваться к перегрузке
	\item Контролируемое увеличение отклонений задач
	\item Балансирование нагрузки между узлами
	\item Стабильность платежей (не скачут резко)
\end{itemize}

\subsection{Сценарий 3: Разнородные ресурсы (Heterogeneous Resources)}

\textbf{Описание}: В этом случае периферийные узлы обладают различной вычислительной мощностью и объёмом памяти. Задачи требуют различное количество ресурсов. 

\textbf{Цель}: Проверить, насколько корректно методы перераспределяют нагрузку в неоднородной инфраструктуре.

\textbf{Ожидаемые результаты}:
\begin{itemize}
	\item Распределение: предпочтение быстрых узлов для срочных задач
	\item Справедливость: компенсация за ограниченные ресурсы
	\item Эффективность: использование разнообразных ресурсов
\end{itemize}

\subsection{Сценарий 4: Отказы узлов (Node Failures)}

\textbf{Описание}: На сороковом шаге эпизода из строя выводится $20\%$ узлов, после чего через $25$ шагов они восстанавливаются.

\textbf{Цель}: Оценить устойчивость алгоритмов к структурным нарушениям среды.

\textbf{Ожидаемые результаты}:
\begin{itemize}
	\item Быстрая адаптация к отказу узла ($< 100$ шагов)
	\item Восстановление производительности
	\item Отсутствие потери справедливости
\end{itemize}

